\section[相对 Fargues--Fontaine 曲线]{相对 Fargues--Fontaine 曲线（11 月 13 日）}

\srcpage{22}

\subsection{回顾上次}

回忆我们的目标：分类「那条」Fargues--Fontaine 曲线 $X_{C,E}$ 上的向量丛。上次我们开始给出线丛的分类，以及斜率 $0$ 的半稳定向量丛的分类。这依赖于对 $\mathcal E\in\mathrm{VB}(X_{C,E})$ 的 $H^0(X_{C,E},\mathcal E)$ 赋予一个「几何结构」；对 $H^1(X_{C,E},\mathcal E)$ 类似。它们称为「Banach--Colmez 空间」。

一个中间目标对后续的一切都至关重要：解释这些「几何结构」。它们称为菱形（diamond）\yizhu{「菱形」是本译本对 diamond 的定译；学界亦作「钻石」，本书统一作「菱形」。}{I（本库 SY02-41 术语表；术语裁定.md）}，同时要对完满空间（perfectoid space）$S$ 定义相对 Fargues--Fontaine 曲线 $X_{S,E}$。这就是我们前进的方向。

参考：[S17]。

\subsection{相对 Fargues--Fontaine 曲线}

回忆：$\Fp$ 上的完满代数就是 $\Fp$ 上的完满 Tate 代数 $R$。因此它有一个拓扑幂零的单位 $\varpi$。（在特征 $p$ 的情形，这恰说明 $R$ 是完满的。）

$\Fp$ 上的完满空间是指 $\Fp$ 上的 adic 空间 $X$，它被开子集 $U=\Spa(R,R^+)\subset X$ 覆盖，且这些 $R$ 完满。

若 $S=\Spa(R,R^+)$ 仿射完满，我们可以模仿 Fargues--Fontaine 曲线的构造，把 $\oo_C$ 换成 $R^+$。确切地说，先作 $\Spa\,W_{\oo_E}(R^+)$，然后看开子空间
\[
Y_{(R,R^+),E}:=\{[\varpi]\neq 0,\ \pi\neq 0\}\subset\Spa\,W_{\oo_E}(R^+).
\]
（回忆 $\pi\in\oo_E$ 是一致化子。）这是 $E$ 上的 adic 空间。

我们偶尔也考虑
\[
\mathcal Y_{(R,R^+),E}:=\{[\varpi]\neq 0\},
\]
它是 $\oo_E$ 上的 adic 空间。于是有包含
\[
\Spa\,W_{\oo_E}(R^+)\supset\mathcal Y_{(R,R^+),E}\supset Y_{(R,R^+),E}.
\]
$\mathcal Y_{(R,R^+),E}$ 与 $Y_{(R,R^+),E}$ 都是解析 adic 空间，意即它们局部是某个 Tate 环的 adic 谱。

存在半径函数（radius function）
\[
\mathrm{rad}\colon\mathcal Y_{(R,R^+),E}\to[0,\infty),\qquad
y\mapsto\frac{\log|\varpi(\tilde y)|}{\log|\pi(\tilde y)|},
\]
其中 $\tilde y$ 是 $y$ 的极大秩 $1$ 泛化\footnote{「秩 $1$」指其值群可嵌入 $\mathbb R$。为取对数我们选取这样的一个嵌入；它当然不唯一，但比值是良定义的。}，并且它把 $Y_{(R,R^+),E}$ 送到 $(0,\infty)$。存在自同态 $\phi_{R^+}$，它自由且完全间断地作用在 $Y_{(R,R^+),E}$ 上：把半径乘以 $q$ 即可看出，
\[
\mathrm{rad}\circ\phi_{R^+}=q\circ\mathrm{rad}.
\]
因此可以取 $Y_{(R,R^+),E}$ 关于 $\phi_R^{\mathbf Z}$ 的商。

\begin{definition}
定义 $X_{(R,R^+),E}=Y_{(R,R^+),E}/\phi_{R^+}^{\mathbf Z}$。它是 $E$ 上的 adic 空间，称为「相对 Fargues--Fontaine 曲线」。
\end{definition}

\begin{example}[等特征]
$E=\Fq((t))$。则 $W_{\oo_E}(R^+)=R^+[[t]]$。所以
\[
\mathcal Y_{(R,R^+),E}=\Spa(R,R^+)\times_{\Spa\,\Fq}\Spa\,\Fq[[t]]=\mathbb D_{\Spa(R,R^+)},
\]
即开单位圆盘。（注意它并非拟紧，这说明仿射空间的纤维积未必仿射。）

子空间 $Y_{(R,R^+),E}=\mathbb D^*_{\Spa(R,R^+)}$ 是去心开单位圆盘。
\end{example}

\srcpage{23}

注意在此情形，$Y_{(R,R^+),E}$ 与 $\mathcal Y_{(R,R^+),E}$ 都有到 $\Spa(R,R^+)$ 的结构映射；但取 $\phi_{R^+}^{\mathbf Z}$ 的商之后就没有了。所以即使在等特征情形，$X_{S,E}$ 也不映射到 $S$。

现在来粘接这一构造。我们断言：对一般的完满空间，图表
\[
\mathcal Y_{S,E}\supset Y_{S,E}\to X_{S,E}=Y_{S,E}/\phi_S^{\mathbf Z}
\]
可以由仿射情形粘接得到。

在正特征 $E=\Fq((t))$ 下这容易看出，此时整张图表都粘接：
\[
\begin{tikzcd}[column sep=large,row sep=small]
\mathcal Y_{S,E} \arrow[r, "="] & S\times_{\Spa\,\Fq}\Spa\,\Fq[[t]] \\
Y_{S,E} \arrow[r, "="] \arrow[u, hook] \arrow[d, two heads] & S\times_{\Spa\,\Fq}\Spa\,\Fq((t)) \arrow[u, hook] \arrow[d, two heads] \\
X_{S,E} \arrow[r, "="] & Y_{S,E}/\phi_S^{\mathbf Z}
\end{tikzcd}
\]

注意 $\Phi_S$ 在拓扑空间 $|S|$ 上是恒等的。所以映射 $|Y_{S,E}|\to|S|$ 下降为拓扑空间的映射 $|X_{S,E}|\to|S|$。在 $p$ 进情形我们想类似地论证。

\begin{remark}
一般地，有「菱形方程」$\mathcal Y^\diam_{S,E}=S\times(\Spa\,\oo_E)^\diam$，且类似的菱形图表可以粘接：
\[
\begin{tikzcd}[column sep=large,row sep=small]
\mathcal Y^\diam_{S,E} \arrow[r, "="] & S\times(\Spa\,\oo_E)^\diam \\
Y^\diam_{S,E} \arrow[r, "="] \arrow[u, hook] \arrow[d, two heads] & S\times(\Spa\,E)^\diam \arrow[u, hook] \arrow[d, two heads] \\
X^\diam_{S,E} \arrow[r, "="] & Y^\diam_{S,E}/\phi_S^{\mathbf Z}
\end{tikzcd}
\]
\end{remark}

\subsection{菱形}

想法是：菱形是（特征 $p$ 的）完满空间按 pro-étale 等价关系取的商。

将有一个从 $\Spa\,\Zp$ 上解析 adic 空间到菱形的函子，记作 $X\mapsto X^\diam$；应当把它看作倾斜函子的推广。回忆倾斜把特征 $0$ 的完满空间送到特征 $p$ 的完满空间。在某种意义上，它像是「忘掉到 $\Qp$ 的结构映射」。

构造 $X^\diam$ 的想法是：取 pro-étale 覆盖 $\tilde X\to X$，其中 $\tilde X$ 完满；然后定义 $X=\tilde X/R$，其中 $R\subseteq\tilde X\times\tilde X$ 完满，且是 pro-étale 等价关系。再定义 $X^\diam=\tilde X^\diam/R^\diam$。

要点在于：倾斜之后 $\mathcal Y_{S,E}$ 会像注 5.3 那样分解为乘积。这与我们对 Fargues--Fontaine 曲线「应当」是什么的直觉相符（参看等特征情形）。

\subsection{完满空间的 pro-étale 映射}

\begin{definition}
设 $f\colon Y\to X$ 是完满空间的映射（可以是混合特征的）。称 $f$ 为有限 étale 的，若对任意开仿射完满子集\footnote{目前不知道完满空间的任一仿射子集是否由完满环表示（定义只要求有这样一个覆盖）。} $U=\Spa(R,R^+)\subset X$（等价地，对这样一个覆盖），原像 $V=f^{-1}(U)=\Spa(S,S^+)\subseteq Y$ 仿射完满，$S$ 是有限 étale $R$-代数，且 $S^+\subset S$ 是 $R^+$ 的整闭包。
\end{definition}

\srcpage{24}
\begin{center}
\figP{024}
\end{center}

为使这一概念行为良好，我们隐含地用了：完满代数上的任意有限 étale 代数自动完满；这是 Faltings 几乎纯度定理的一个版本。（在特征 $p$ 下这显然。）所以 $\Spa(R,R^+)_{\mathrm{fét}}\cong(\Spec R)_{\mathrm{fét}}$。

\begin{definition}
设 $f\colon Y\to X$ 是完满空间的映射（可以是混合特征的）。称 $f$ 为 étale 的，若它在 $X$ 与 $Y$ 上局部地形如
\[
\begin{tikzcd}[column sep=large,row sep=large]
Y \arrow[r, hook, "\text{开浸入}"] & Y' \arrow[dl, "\text{有限 étale}"'] \\
U \arrow[r, hook, "\text{开浸入}"'] & X
\end{tikzcd}
\]
\end{definition}

étale 映射的复合仍是 étale 的，这并不平凡但成立。

\begin{warning}
注意：对概形，类似的刻画是错的。考虑曲线的分歧映射，再在分歧点处打洞。定义域的任一 Zariski 开子集都能看到分歧点的邻域，解析开子集则不然。
\end{warning}

\begin{definition}
设 $f\colon Y\to X$ 是完满空间的映射（可以是混合特征的）。称 $f$ 为 pro-étale 的，若它在 $X$ 与 $Y$ 上局部是仿射 pro-étale 的；即 $Y=\Spa(S,S^+)$ 是形如
\[
\varprojlim_i\Spa(S_i,S_i^+)\to X=\Spa(R,R^+)
\]
的余滤过极限，且对每个 $i$，$f_i\colon(\Spa,S_i,S_i^+)\to\Spa(R,R^+)$ 是 étale 的。\yizhu{原文此处印作 $(\Spa,S_i,S_i^+)$，按上下文应为 $\Spa(S_i,S_i^+)$；译文照原文保留。}{A（原文 §5.4 p.24）}
\end{definition}

我们只限于仿射的情形，因为它们有一切（连通）极限；特别地 $S^+=(\varinjlim S_i^+)^\wedge_\varpi$，且 $S=S^+[1/\varpi]$。

\begin{example}
设 $p\neq 2$。令 $X=\Spa\,C\langle T^{1/p^\infty}\rangle$，$Y=\Spa\,C\langle T^{1/2p^\infty}\rangle$。映射 $Y\to X$ 看起来在原点处分歧。然而我们将造出 pro-étale 覆盖 $\tilde X\to X$，使 $Y\to X$ 的基变换是仿射 pro-étale 的。这说明：虽然 $Y\to X$ 不是 pro-étale 的，它在 pro-étale 拓扑下局部是 pro-étale 的；这很奇怪，与概形的通常 étale 拓扑大不相同。
\[
\begin{tikzcd}[column sep=large,row sep=large]
\tilde Y \arrow[r, "\text{pro-étale}"] \arrow[d, "\text{仿射 pro-étale}"'] & Y \arrow[d] \\
\tilde X \arrow[r, "\text{pro-étale}"'] & X
\end{tikzcd}
\]
令 $U_n=\{|T|\le 1/p^n\}$，$U_{n,n+1}=\left\{\frac{1}{p^{n+1}}\le|T|\le\frac{1}{p^n}\right\}\subset U_n$。这些都是有理子集，因而仿射完满。对每个 $n$，令 $X_n=U_{0,1}\sqcup U_{1,2}\sqcup\cdots\sqcup U_{n-1,n}\sqcup U_n$。每个 $X_n\to X$ 都是 étale 覆盖。
\end{example}

\srcpage{25}

现在考虑 $\tilde X=\varprojlim_n X_n\to X$。注意 $\pi_0\tilde X=\mathbb N\cup\{\infty\}$。$\tilde X$ 在 $n\in\mathbb N$ 上的纤维是 $U_{n,n+1}$，在 $\infty$ 上的纤维是 $\Spa\,C$，即原点。所以 $Y\times_X\tilde X\to Y$ 是仿射 pro-étale 的。

\begin{example}
任一 Zariski 闭浸入都是仿射 pro-étale 的。设 $X=\Spa(R,R^+)\supset V(f)=Z=\Spa(S,S^+)$。令 $S=R/(f^{1/p^\infty})\supseteq S^+$ 为 $R^+$ 的整闭包。令 $U_n=\{|f|\le 1/p^n\}\subset X$ 为有理开子集。考虑 $U_\infty:=\bigcap_n U_n=\varprojlim_n U_n\subset X$。在 $U_\infty$ 上处处 $|f|=0$，故 $f=0$（完满空间既约，甚至是 uniform 的）。所以 $U_\infty=V(f)$。
\end{example}

换一种稍有不同的说法：在 $U_\infty$ 上，函数 $f$ 越来越能被 $\varpi$ 整除，于是在完备化中为 $0$。

注意：代数几何中类似的构造会给出 Hensel 化。

\begin{example}
$f\colon\Spa\,C\langle T^{1/p^\infty}\rangle\to C$ 不是 pro-étale 的，这是维数原因。（pro-étale 映射的纤维是射有限集。）\yizhu{原文把目标印作 $C$；按上下文应为 $\Spa\,C$。译文照原文保留。}{A（原文 §5 p.25）}
\end{example}

\begin{theorem}
仿射完满空间的映射 $f\colon Y\to X$ 在 $X$ 上 pro-étale 局部是仿射 pro-étale 的，当且仅当：
\begin{itemize}
\item 对所有几何（秩 $1$）点 $\Spa(C,\oo_C)\to X$，纤维积
\[
Y\times_X\Spa(C,\oo_C)\to\Spa(C,\oo_C)
\]
是仿射 pro-étale 的；等价地，同构于 $\mathcal S\times\Spa(C,\oo_C)=\varprojlim_i(\mathcal S_i\times\Spa(C,\oo_C))$，其中 $\mathcal S=\varprojlim_i\mathcal S_i$ 是射有限集。
\end{itemize}
\end{theorem}

\begin{definition}
这样的映射称为 quasi-pro-étale 的。要点是：这一点可以在 pro-étale 拓扑下局部地检查。
\end{definition}

\subsection{v-拓扑}

\begin{definition}
映射 $f\colon Y\to X$ 称为 v-覆盖，若对任意仿射 $U\subset X$，存在拟紧开子集 $V\subset Y$ 使 $|V|\to|U|$ 满射。若 $f$ 是 v-覆盖且 quasi-pro-étale，则称 $f$ 为 pro-étale 覆盖。
\end{definition}

\begin{theorem}
\begin{enumerate}[label=(\arabic*)]
\item 预层 $X\mapsto\oo_X(X)$ 与 $X\mapsto\oo_X^+(X)$ 是完满空间上 v-拓扑的层。
\item 对任意完满空间 $X$，$\Hom(-,X)$ 是 v-拓扑的层。
\item $X\mapsto\mathrm{VB}(X)$ 是 v-叠。
\item 若 $X=\Spa(R,R^+)$ 仿射，则
\[
H^i_v(X,\oo_X)=0\quad(i>0),\qquad H^i_v(X,\oo_X^+)\overset{a}{=}0\quad(i>0).
\]
\end{enumerate}
这里 $\overset{a}{=}0$ 表示「几乎为 $0$」，即被所有 $\varpi^{1/p^n}$ 杀死。
\end{theorem}

\begin{remark}
v-拓扑是 fpqc 拓扑的类比，后者通常包含平坦性条件与拓扑条件。注意在完满世界中「一切都是平坦的」，所以只剩下拓扑条件。

第 (2) 部分暗示：v-拓扑本质上是典范拓扑。在仿射完满空间上，确实有：v-拓扑就是典范拓扑（使可表对象为层的最细拓扑）。
\end{remark}

\srcpage{26}

\begin{proof}[定理 5.14 的证明概要]
首先证明 $X\mapsto\oo_X^+(X)$ 对 étale 拓扑是层，且当 $X$ 仿射、$i>0$ 时 $H^i_{\mathrm{Ét}}(X,\oo_X^+)$ 几乎为 $0$。为此分开启覆盖与有限 étale 覆盖两种情形，用 de Jon 与 van der Put 的组合论证。开覆盖情形基本上是「经典的」；有限 étale 情形基本上归结为 Faltings 的几乎纯度定理。

然后对 $\oo_X^+/\varpi$ 得到类似断言。借助滤过余极限，它推广到仿射 pro-étale 的情形（$\oo_X^+/\varpi$ 是离散的，与余极限交换得好）。所以 $\oo_X^+/\varpi$ 作为 pro-étale 层性质良好。再提升到 $\oo_X^+=\varprojlim_n\oo_X^+/\varpi^n$，进而到 $\oo_X=\oo_X^+[1/\varpi]$。

接着利用：pro-étale 局部地，v-覆盖在 $\oo^+/\varpi$ 层面是忠实平坦的。这看起来令人惊讶，因为我们没有假设任何平坦性。要点在于：在 pro-étale 局部，你把空间打碎到基本上只是一堆点，而那里平坦性是自动的。所以「pro-étale 局部地一切都是平坦的」。然后化归到通常的忠实平坦下降理论。
\end{proof}
